\section{Multi-tenancy in Vector Databases}
\label{sec:multitenancy}

In production, vector database tenants rarely live in isolation. A
typical deployment maintains a shared corpus of common knowledge,
for example a public document collection or a base code repository,
that every tenant queries. On top of that substrate, each tenant
introduces a private overlay. Some tenants override specific
records by substituting a tenant-local version of a shared vector
under the same logical key, while others append entirely new
documents that exist only for that tenant. We refer to the first
pattern as per-tenant versioning and to the second as forking. In
both cases a query issued by a tenant must be answered against the
union of shared and tenant-private vectors, so the storage layer
cannot simply isolate tenants into disjoint collections without
either losing the substrate or duplicating it across all tenants.

This setting raises two design questions. First, at what
granularity should tenant data be separated inside the database,
and what does each granularity choice imply for storage and query
cost? Second, how should the shared corpus and the per-tenant
overlay be represented so that the shared portion is paid for once
while keeping query latency predictable? The remainder of this
section reviews the data granularity primitives Milvus exposes.
Section~\ref{sec:setup} then formalizes the two layout options within this design space and the experimental setup we use to
characterize each one.

\subsection{Data layout for multi-tenant collections}

\mypar{Collection, partition, and partition key.}
Milvus~\cite{wang2021milvus} exposes three nested units of
organization for tenant data. A \emph{collection} is the heaviest
unit, carrying its own schema, index configuration, and resource
allocation, so two tenants placed in two different collections share
nothing at the storage layer. A \emph{partition} is a sub-grouping
inside a single collection, and a query can be constrained to a list
of partitions so that only a subset of the data is consulted. The
\emph{partition key} is a lightweight variant of the partition
mechanism in which a scalar field, typically a tenant identifier, is
hashed into one of a fixed number of auto-managed partitions and
routes vectors to partitions automatically. The partition key is the
mechanism Milvus recommends for deployments with many tenants
because it avoids the operational overhead of explicitly managing
per-tenant partitions.

\mypar{Segments and how a collection lives on disk.} Below the
partition level sits the \emph{segment}, the physical storage unit
introduced in Section~\ref{sec:background}. Each segment holds a
single data file together with its own vector index, and segments
never cross partition boundaries. A partition therefore owns a set
of segments, and a collection is the disjoint union of those
per-partition segment sets. A partition key does not change this
physical picture, it only shifts the routing decision into the
database so that every vector is hashed to one of the auto-managed
partitions before being assigned to a segment. As a result, the
choice of granularity, whether collection, explicit partition, or
partition key, controls how much storage and index state two tenants
share, while the underlying segment, data file, and index structure
remain identical.

\subsection{Representing shared and per-tenant data}

Given the granularity primitives reviewed in
Section~\ref{sec:multitenancy}, there are two natural ways to
represent a shared corpus together with per-tenant overlays. The
first, \emph{fully replicated}, gives each tenant its own collection
that contains both the shared vectors and the tenant's private vectors.
This layout offers strong isolation, since two tenants share no
schema, index, or in-memory state, but pays a storage cost
proportional to the number of tenants because the shared corpus is
duplicated in every collection. The second, \emph{shared partition},
keeps a single collection in which the shared vectors live in one
partition and each tenant's private vectors live in a dedicated
partition. The shared corpus is then stored, indexed, and cached
only once, but the partitions sit inside the same collection and
contend for the same memory and page cache, so cross-tenant
interference at query time is not bounded by construction.

These two options frame the experiments in the remainder of this
report. Under the fully replicated layout, sharing storage across
tenants is only possible if the segments produced by independent
ingest streams of near-identical data actually align, so we measure
how much segment-level overlap stock Milvus exhibits in this regime.
Under the shared partition layout, the shared corpus is stored only
once but tenants contend for the same memory, so we measure how
query throughput behaves once two tenants query concurrently against
a single collection. The next two subsections describe the workloads
and configurations used for each of these experiments.